% tab:sources
\begin{table}[htbp]
\centering\footnotesize
\caption[Volume indicators and sources by sector]{\textsc{Volume indicators and sources by sector}}
\label{tab:sources}
\begin{threeparttable}
\begin{tabular}{>{\RaggedRight}p{2.6cm}>{\RaggedRight}p{5.3cm}>{\RaggedRight}p{5.3cm}c}
\toprule
Sector & Interwar series, 1919--40 (1938 weights) & Soviet-period series, 1937--95 (1990 weights) & App. \\
\midrule
Crops & Physical output of 15 crops weighted by 1938 value-added (Valge; 1939 agricultural survey); yields from \cite{klesment2007} & Physical output; weights from World Bank (1996) and Eurostat agricultural accounts; 1945--65 minor crops interpolated between quinquennial observations & A.1 \\
Livestock & Physical output of meat, milk, eggs and wool; 1938 value-added weights & Physical output; during 1941--44 animal numbers multiplied by pre-war yields & A.1 \\
Forestry and fishing & Timber removals and fish catch; timber for 1919--21 and fish for 1919--27 extrapolated from the crops-and-livestock index & Timber removals and fish catch; 1941--44 interpolated & A.1 \\
Industry & Employment adjusted for real wages (no usable annual output data; Valge's current-price series cannot be deflated by branch) & 101 physical product series in 11 branches; branch weights from the 1990 benchmark, product weights from 1986 archival output data (ERA) & A.2 \\
Construction & Domestic cement use (production plus net imports), as in Valge; 1919--21 from population & Weighted index of material inputs (1987 Soviet input--output weights), following the CIA procedure & A.3 \\
Transport and communications & Railway passenger- and tonne-kilometres, merchant fleet tonnage, postal items, telegrams and telephone calls & Freight tonne-kilometres and passenger-kilometres for all modes; postal, telephone and broadcasting indicators; 1941--45 interpolated between 1940 and the first post-war observation, with 1943--45 movements from Lithuanian ratios (Section~\ref{sec:method}) & A.4 \\
Trade, finance and insurance & Marketed value-added of agriculture and industry plus deflated imports (`commerce' index); finance assumed to move with commerce & Soviet retail and wholesale turnover at comparable prices (archival and published); finance and insurance from employment with a productivity adjustment & A.5 \\
Housing & Valge's housing value-added for 1929--38, extrapolated with the urban population & Housing stock from construction and population data & A.5 \\
Public administration, defence, health and education & Civil services assumed to move with GDP; defence from military personnel (15,900 in 1921--39) & Employment by branch (Soviet yearbooks, archival records); defence from the 1940 and 1990 totals & A.6 \\
Other services & Total population & Employment & A.6 \\
\bottomrule
\end{tabular}
\begin{tablenotes}[para,flushleft]
\footnotesize
\textit{Sources:} Interwar: \cite{valge2003all}; \cite{klesment2007}; interwar statistical yearbooks and agricultural surveys. Soviet period: \textit{Eesti NSV rahvamajandus} statistical yearbooks (1957--89), \textit{Eesti statistika aastaraamat} (1990--98), Estonian National Archives (ERA, fond R-10 and R-1969), German occupation statistical bulletins (\textit{Statistische Berichte f\"ur das Ostland}, 1942--44), and World Bank (1993, 1996). Full source lists by series are given in the online appendix, sections A.1--A.6.
\end{tablenotes}
\end{threeparttable}
\end{table}
